\slajdid{08-s064}
\begin{frame}[label=08-s064]{Induktivnost}
\gusttekst
\begin{columns}[c,onlytextwidth]\column{.23\textwidth}$\displaystyle u_L(t)=L\frac{di_L(t)}{dt}$
\column{.40\textwidth}\centering\begin{tikzpicture}[x=1cm,y=1cm,font=\sitneoznake,>=Latex]
\draw(0,0)circle(.25);\node at(.12,.4){$+$};\node[left]at(-.3,0){$u_i(t)$};\draw(0,.25)--(0,.85)--(1,.85);
\foreach\x in{1,1.25,1.5,1.75}{\draw(\x,.85)..controls(\x,.99)and(\x+.25,.99)..(\x+.25,.85);}
\draw(2,.85)--(3.6,.85);\node[above]at(1.5,1){$L$};\draw[->](.3,1)--(.75,1)node[midway,above]{$i_L(t)$};
\draw[<->](1,.55)--(2,.55)node[midway,below]{$u_L(t)$};\node[left]at(1,.55){$+$};
\draw(0,-.25)--(0,-.85)--(3.6,-.85);\draw(2.9,.85)--(2.9,.5)--(3.02,.4)--(2.78,.24)--(3.02,.08)--(2.78,-.08)--(3.02,-.24)--(2.78,-.4)--(2.9,-.5)--(2.9,-.85);\node[left]at(2.7,0){$R$};
\draw[<->](3.6,.75)--(3.6,-.75);\draw[fill=white](3.6,.85)circle(.06);\draw[fill=white](3.6,-.85)circle(.06);\node[right]at(3.7,.6){$+$};\node[right]at(3.7,0){$u_o(t)$};
\end{tikzpicture}

\column{.34\textwidth}\centering\begin{tikzpicture}[x=.65cm,y=.65cm,font=\sitneoznake,>=Latex]
\draw[->](-1.5,0)--(4.6,0)node[below]{$t$};\draw[->](0,-.15)--(0,2.5)node[left]{$u_i(t)$};\draw(-1,.3)--(.7,.3)--(.7,2)--(3.9,2);\draw[dashed](.7,-.15)--(.7,.3);\node[below]at(.7,-.15){$t_0$};\node[above]at(2.3,2){$=V_{OH}$};\node[above]at(-.6,.3){$=V_{OL}$};
\end{tikzpicture}

\end{columns}\medskip
Da bi se struja kroz induktivnost promenila potreban je beskonačno veliki napon,\par\medskip
Ako postoji napon na kalemu biće i promene struje kroz njega.\par\medskip
Isto kao i u slučaju kapacitivnosti. Ako je vreme pre trenutka promene dovoljno dugo trajalo (beskonačno), kolo će ući u stacionarno stanje u kome nema više promena i karakterisano je
\[u_L(t)=0,\qquad i_L(t)=const\]
\end{frame}
